\documentclass[11pt]{article}
\usepackage[T1]{fontenc}
\usepackage[cmex10]{amsmath}
\usepackage{amssymb,amsthm,booktabs}
\usepackage[margin=1in]{geometry}
\usepackage[hidelinks]{hyperref}
\DeclareFontShape{OMX}{cmex}{m}{n}{<->cmex10}{}
\newtheorem{theorem}{Theorem}
\newtheorem{lemma}{Lemma}
\hypersetup{pdftitle={Real-rootedness for balanced arrays with three columns and row cap two},pdfauthor={ProgramComputer}}
\setlength{\emergencystretch}{1em}
\widowpenalty=4000
\clubpenalty=4000
\title{Real-rootedness for balanced arrays\\ with three columns and row cap two}
\author{ProgramComputer}
\date{November 25, 2026}

\begin{document}
\maketitle

\begin{abstract}
For every \(n\), the polynomial that counts binary \(n\times3\) arrays with a common column sum and row sums at most two has only simple negative zeros. This is the \(U(2,3)\) specialization of Alexandersson's balanced-array conjecture. Dividing the coefficient at column sum \(k\) by \(\binom nk\) yields an algebraic generating function with four branch points. Exact contour bounds and a Legendre differential-equation estimate give the full number of sign changes for every \(n\ge20\). A finite multiplier restores the original coefficients, and exact rational Sturm calculations verify the twenty remaining rows.
\end{abstract}

\section{Statement}

For \(n,k\in\mathbb Z_{\ge0}\), let \(a_{n,k}\) count the \(n\times3\) zero--one matrices in which every column has sum \(k\) and every row has sum at most two. Rows and columns are indexed. Allow the empty array at \(n=0\), and set
\[
P_n(q)=\sum_{k\ge0}a_{n,k}q^k
 =\sum_{k\ge0}[x_1^kx_2^kx_3^k]
 \bigl(1+e_1(x_1,x_2,x_3)+e_2(x_1,x_2,x_3)\bigr)^n q^k,
\]
where \(e_j\) is the \(j\)th elementary symmetric polynomial. The permitted row supports are the independent sets of the uniform matroid \(U(2,3)\). For \(0\le k\le n\), set \(g_{n,k}=a_{n,k}/\binom nk\) and \(G_n(q)=\sum_{k=0}^n g_{n,k}q^k\).

\begin{theorem}\label{thm:main}
For every \(n\ge0\), \(P_n\) and \(G_n\) have degree \(D_n=\lfloor 2n/3\rfloor\) and each has \(D_n\) distinct negative zeros. The statement includes the constant rows \(n=0,1\).
\end{theorem}

Alexandersson conjectures real-rootedness for all row caps and all column counts, and identifies \(U(2,4)\) as the first target in the original issue~\cite{original,catalog}. Here we establish the three-column specialization for every row count. Consecutive-row interlacing is a separate question.

\section{The normalized generating function}\label{sec:generating}

Fixing the third column to a particular $k$-subset and excluding rows containing three ones by inclusion-exclusion gives
\begin{equation}
g_{n,k}=\sum_{j=0}^k(-1)^j\binom kj\binom{n-j}{k-j}^2.
\label{eq:normalized-count}
\end{equation}
The degree of both polynomials is $D_n=\lfloor2n/3\rfloor$. The bound follows by counting ones. For $n>0$ and $3k\leq2n$, list the $3k$ positions $0,1,\ldots,3k-1$ modulo $n$ and give successive groups of $k$ to the three columns. Each column receives distinct row labels since $k\le n$, and each row occurs at most twice. Hence every coefficient up to $D_n$ is positive.

Put $u=n-k$ in \eqref{eq:normalized-count} and sum the binomial transform to obtain
\[
\sum_{k\ge0}g_{u+k,k}t^k
 =\frac1{1+t}\sum_{\ell\ge0}\binom{u+\ell}{u}^2
   \left(\frac{t}{1+t}\right)^\ell
 =(1+t)^u L_u(1+2t),
\]
where $L_u$ is the Legendre polynomial. The last equality follows from
\begin{align*}
\sum_{\ell\ge0}\binom{u+\ell}{u}^{2}z^\ell
   &=(1-z)^{-2u-1}\sum_{i=0}^{u}\binom ui^2z^i,\\
\sum_{i=0}^{u}\binom ui^2z^i
   &=(1-z)^uL_u\left(\frac{1+z}{1-z}\right).
\end{align*}
Set $t=xq$, multiply by $x^u$, and sum over $u$ by the Legendre generating function. This gives
\begin{equation}
\sum_{n\ge0}G_n(q)x^n=D(x,q)^{-1/2},\qquad
D(x,q)=1-2x+(1-6q)x^2+(2q-4q^2)x^3+q^2x^4.
\label{eq:gf}
\end{equation}
The branch equals $1$ at $x=0$. See \cite[(15.8.1), (18.5.8), and (18.12.11)]{dlmf} for these identities.

\section{Four reciprocal branch points}\label{sec:branch}

Throughout the proof let $q=-a<0$. The factorization
\[
D(s^2,q)=\big((1+s)^2+qs^3(s+2)\big)\big((1-s)^2+qs^3(s-2)\big)
\]
shows that $\lambda=(v-1)^2/4$ satisfies $D(1/\lambda,-a)=0$ whenever $(v^2-1)^2-16av=0$. Indeed, substituting $s=2/(v-1)$ in the first factor gives $((v^2-1)^2-16av)/(v-1)^4$. For each $a>0$ there is a unique $t\in(0,1)$ such that
\[
a=\frac{(t+t^{-1})(t^{-2}-t^2)}{16},\qquad
S=t+t^{-1}.
\]
Both factors in the numerator are positive and strictly decreasing on $(0,1)$, and their product decreases from $+\infty$ to $0$. For this $t$ the polynomial in $v$ factors as
\[
(v^2-Sv+t^2)(v^2+Sv+t^{-2}).
\]
The first quadratic has distinct positive roots $r<1<R$, since it is negative at $1$, positive at $0$, and positive for large positive $v$. The second has a nonreal conjugate pair $w,\bar w$, since $4t^{-2}-S^2=t^{-2}(1-t^2)(3+t^2)>0$. Take $w$ in the lower half-plane. Define
\begin{align*}
\lambda_1&=(R-1)^2/4, & \lambda_2&=(1-r)^2/4,\\
\lambda_+&=(w-1)^2/4=\rho e^{i\phi}, & \lambda_-&=\overline{\lambda_+}.
\end{align*}
Set $\tau=\lambda_1/\rho$ and $\sigma=\lambda_2/\rho$.
Here $\rho=(1+S+t^{-2})/4>1$ because $\Re w=-S/2$ and $|w|^2=t^{-2}$. We have $\lambda_1>\lambda_2>0$ since $R+r=S>2$. Also
\begin{equation}
4(\rho-\lambda_1-\lambda_2)=t^2+3S-3>0,\qquad 0<\sigma<1/4.
\label{eq:modulus-order}
\end{equation}
For the first identity, use $R^2+r^2=S^2-2t^2$ and $S^2=t^2+t^{-2}+2$. For the second inequality, $\lambda_2<1/4$. The argument of $\lambda_+$ is $2\arctan(|\Im w|/(1-\Re w))$. More explicitly, with $m=(1-t)/(1+t)$,
\begin{equation}
\tan^2(\phi/2)=\frac{(1-t)(3+t^2)}{(1+t)^3}=m+m^2+m^3,\qquad
\rho=\frac{1+m^2}{(1+m)(1-m)^2}.
\label{eq:param}
\end{equation}
In particular $\phi$ increases strictly from $0$ to $2\pi/3$ as $a$ increases from $0$ to $+\infty$. These four reciprocal branch points are distinct, so
$D(x,-a)=\prod_{\nu\in\{+,-,1,2\}}(1-\lambda_\nu x)$.
The ratios satisfy
\begin{equation}
\tau+\sigma=\frac2{\rho}-2\cos\phi
  =\frac{2m(1-m+m^2+m^3)}{(1+m)(1+m^2)}\leq 2m,
\qquad
\lim_{\phi\to 2\pi/3}(\tau,\sigma)=(1,0).
\label{eq:ratios}
\end{equation}
The sum formula uses the coefficient of $x$ in \eqref{eq:gf}. At the upper endpoint $t\to0$, one has $tr\to0$, $tR=1+t^2-tr\to1$, and $4t^2\rho\to1$, proving the ratio limit.

Use principal square roots in the amplitude
\[
A=\prod_{\nu\ne+}(1-\lambda_\nu/\lambda_+)^{-1/2}.
\]
Set
\begin{equation}
\eta=\tfrac12\big(\arg(1-\tau e^{-i\phi})+\arg(1-\sigma e^{-i\phi})\big),
\qquad
\Psi_n(\phi)=(n+\tfrac12)\phi-\tfrac{\pi}{4}-\eta.
\label{eq:phase}
\end{equation}
The two arguments in $\eta$ are positive and less than $\pi/2$ by \eqref{eq:modulus-order}. Also $1-e^{-2i\phi}=2\sin\phi\,e^{i(\pi/2-\phi)}$. Thus $e^{in\phi}A=|A|e^{i\Psi_n(\phi)}$. The continuous extension to the closed phase interval has values
\begin{equation}
\Psi_n(0)=-\pi/4,\qquad \Psi_n(2\pi/3)=2n\pi/3.
\label{eq:phase-ends}
\end{equation}
For the first limit, \eqref{eq:ratios} gives $\tau,\sigma\to0$. For the second, \eqref{eq:ratios} gives $\eta\to\arg(1-e^{-2\pi i/3})/2=\pi/12$.

\section{Signs at the phase nodes}\label{sec:signs}

\begin{lemma}\label{lem:sign}
For \(n\ge20\), if \(0<\phi<2\pi/3\) and \(\Psi_n(\phi)=k\pi\) with integer \(k\ge1\), then \(G_n(-a)\) has sign \((-1)^k\).
\end{lemma}

\begin{proof}
We handle the two angular ranges in turn.

\medskip
\noindent\textit{The range \(\pi/6\le\phi<2\pi/3\).}
For integers $j\geq0$ let $c_j=\binom{2j}{j}/4^j$. The beta integral gives
\[
\pi c_n=\int_0^1 u^{n-1/2}(1-u)^{-1/2}\,du.
\]
For real $u\in(0,1]$ let
\[
A_+(u)=\prod_{\nu\ne+}\left(1-\frac{\lambda_\nu}{u\lambda_+}\right)^{-1/2}.
\]
The real branch point contribution will use $A_1(u)$, defined by replacing $+$ with $1$ in this formula for $\kappa<u\leq1$, where $\kappa=\sigma/\tau$. All factors use principal square roots. Deforming the Cauchy coefficient contour to the cuts gives
\begin{equation}
G_n(-a)=
  2\operatorname{Re}\frac{\lambda_+^n}{\pi}
       \int_0^1\frac{u^{n-1/2}}{\sqrt{1-u}} A_+(u)\,du
 +\frac{\lambda_1^n}{\pi}
       \int_\kappa^1\frac{u^{n-1/2}}{\sqrt{1-u}} A_1(u)\,du .
\label{eq:cut}
\end{equation}
Use the product of the four principal reciprocal square roots for the germ in \eqref{eq:gf}. For the conjugate pair use rays $x=y/\lambda_\pm$, $y\geq1$. For the positive real branch points the ray jumps of the two square roots cancel beyond $x=1/\lambda_2$, leaving the interval $[1/\lambda_1,1/\lambda_2]$. On the upper side of a cut in the oriented variable $y=\lambda_\nu x$, the singular factor $(1-y)^{-1/2}$ has value $i(y-1)^{-1/2}$, and the lower value is its negative. The upper side is integrated with $y$ increasing. On the real interval the other factors yield $A_1(1/y)>0$. The small endpoint arcs vanish because the square-root singularities are integrable. On the outer circle of radius $R$, the generating function is $O(R^{-2})$ for fixed $a>0$ and the factor $x^{-n-1}$ makes its coefficient integral vanish as $R\to\infty$. Substituting $u=1/y$ in the jumps gives \eqref{eq:cut}.

For the nonreal integral, a factor $(1-z/u)^{-1/2}$ with fixed $z\ne0$ has logarithmic derivative of absolute value at most $|z|/(2u|u-z|)$. This ratio is at most $1/(2u)$ when $\Re z\leq0$ and at most $1/(2u|\sin\arg z|)$ otherwise. For the two real-root factors of $A_+(u)$ the angle of $z$ is $-\phi$, giving a bound $1/u$ for each. For the conjugate factor the angle is $-2\phi$. Where its real part is positive, the angle bound gives $|\sin(2\phi)|\geq\sqrt3/2$, so the derivative bound is $1/(\sqrt3 u)$ in all cases. Thus
\begin{equation}
\left|\frac{A_+'(u)}{A_+(u)}\right|<\frac3u,\qquad
\frac{|A_+(u)-A|}{|A|}\leq u^{-3}-1.
\label{eq:amp}
\end{equation}
Indeed the first bound gives $|A_+(u)|\leq |A|u^{-3}$, and integrating the resulting bound on $|A_+'|$ gives the second. Dividing by $c_n|\lambda_+|^n|A|$, the error when replacing the first integral with $c_n \lambda_+^n A$, before taking $2\operatorname{Re}$, is at most
\begin{equation}
\frac{n(n-1)(n-2)}{(n-\frac12)(n-\frac32)(n-\frac52)}-1<\frac1{10},
\qquad n\geq20.
\label{eq:beta}
\end{equation}
This is the beta moment of $u^{-3}-1$. Each of the three factors of the first term decreases in $n$, and at $n=20$ their product is $54720/50505<11/10$.

The real contribution is positive and admits a uniform upper bound even when $\tau$ approaches $1$. Its integrand satisfies
\[
A_1(u)=\frac{\tau u}{|1-\tau u e^{-i\phi}|}\sqrt{\frac{u}{u-\kappa}}
 \leq \frac{\tau}{|1-\tau e^{-i\phi}|}\sqrt{\frac{u}{u-\kappa}}.
\]
For the inequality, $ |1-\tau u e^{-i\phi}|^2-u^2|1-\tau e^{-i\phi}|^2
=(1-u)(1+u-2\tau u\cos\phi)\geq0$. The prefactor in this upper bound satisfies
\begin{equation}
\frac{\tau}{|A|\,|1-\tau e^{-i\phi}|}=\tau \sqrt{\frac{2\sin\phi\,|1-\sigma e^{-i\phi}|}{|1-\tau e^{-i\phi}|}}
 \ \leq\ \sqrt{5/2}.
\label{eq:real-prefactor}
\end{equation}
Here $|1-\tau e^{-i\phi}|\geq\sin\phi$, $|1-\sigma e^{-i\phi}|\leq1+\sigma<5/4$, and $\tau<1$. Set $u=\kappa+(1-\kappa)v$ in the remaining integral. With $V$ distributed on $(0,1)$ with density $1/(\pi\sqrt{v(1-v)})$, the real contribution divided by $\rho^n |A|$ is at most
\begin{equation}
\sqrt{5/2}\,\mathbb E[(\sigma+(\tau-\sigma)V)^n]
  \ \leq\ \sqrt{5/2}\,(c_n+2^{-n}).
\label{eq:arcsine}
\end{equation}
For $v\geq1/2$, the linear expression is at most $v$ because $\tau+\sigma<1$. For $v\leq1/2$, it is at most $1/2$ by the same inequality and $\tau>\sigma$. Finally $\mathbb E[V^n]=c_n$.

The elementary bounds $1/(2\sqrt j)\leq c_j\leq1/\sqrt{j+1}$ for $j\geq1$ follow by induction from the quotient $c_{j+1}/c_j=(2j+1)/(2j+2)$. For $n\geq20$ they give $2^{-n}/c_n\leq 2n\,2^{-n}<1/1000$. Therefore, in units of $\rho^n|A|c_n$, the difference between $G_n(-a)$ and its leading oscillation $2\rho^n|A|c_n\cos\Psi_n(\phi)$ is less than
\[
\frac15+\frac85(1+1/1000)<2.
\]
In particular $G_n(-a)$ has sign $(-1)^k$ wherever $\Psi_n(\phi)=k\pi$ in this phase range.

\medskip
\noindent\textit{The range \(0<\phi\le\pi/6\).}
In this interval, \eqref{eq:param} gives $m\leq\tan^2(\pi/12)=7-4\sqrt3<1/12$. By \eqref{eq:ratios}, $\tau,\sigma<1/6$. The generating function \eqref{eq:gf} and the Legendre generating function give the exact convolution
\begin{equation}
\frac{G_n(-a)}{\rho^n}=\sum_{\ell=0}^n h_\ell L_{n-\ell}(\cos\phi),\qquad
\sum_{\ell\geq0}h_\ell z^\ell=(1-\tau z)^{-1/2}(1-\sigma z)^{-1/2}.
\label{eq:legendre-conv}
\end{equation}
Here $h_0=1$ and $0\leq h_\ell\leq6^{-\ell}$, by bounding the two nonnegative binomial series coefficientwise by $(1-z/6)^{-1/2}$. Each argument in \eqref{eq:phase} is at most $\arctan((1/6)/(1-1/6))<1/5$, so $0<\eta<1/5$. At a node $\Psi_n(\phi)=k\pi$ with $k\geq1$, we have
\begin{equation}
(n+\tfrac12)\phi=k\pi+\pi/4+\eta\ \geq\ 5\pi/4.
\label{eq:node-lower-bound}
\end{equation}

For $N\geq1$ set $y_N(\theta)=\sqrt{\sin\theta}\, L_N(\cos\theta)$. The Legendre differential equation~\cite[Table 18.8.1]{dlmf} gives
\[
y_N''+\left((N+\tfrac12)^2+\frac1{4\sin^2\theta}\right)y_N=0.
\]
At $\theta=\pi/2$ the value and derivative agree with those of
$b_N\cos((N+\tfrac12)\theta-\pi/4)$, where
\[
b_N=\begin{cases}
 c_{N/2},&N\text{ even},\\
 \dfrac{N}{N+1/2}\,c_{(N-1)/2},&N\text{ odd}.
\end{cases}
\]
Indeed the Legendre generating function at $0$ gives $L_{2j}(0)=(-1)^jc_j$, and its derivative gives $L'_{2j+1}(0)=(-1)^j(2j+1)c_j$. Variation of constants for the free equation $y''+(N+\tfrac12)^2y=0$ gives a Volterra integral with kernel bounded in absolute value by $1/(N+\tfrac12)$. Integrating backward from $\pi/2$ and applying Grönwall's inequality gives
\begin{equation}
\left|\frac{y_N(\phi)}{b_N}-
       \cos((N+\tfrac12)\phi-\pi/4)\right|
 \ \leq\ \exp\left(\frac{\cot\phi}{4(N+1/2)}\right)-1.
\label{eq:legendre-error}
\end{equation}
Explicitly, $|y_N(\theta)|/b_N\leq1+\int_\theta^{\pi/2} f(s)|y_N(s)|/b_N\,ds$ with $f(s)=1/(4(N+\tfrac12)\sin^2s)$. Inserting the bound $\exp(\int_\theta^{\pi/2}f)$ into the same integral for the difference gives \eqref{eq:legendre-error}.

For $N\geq n/2$, by \eqref{eq:node-lower-bound} and $\cot\phi\leq1/\phi$ the exponent in \eqref{eq:legendre-error} is at most $2/(5\pi)$. Its exponential is less than $7/6$, using $\pi>3$ and $e^{2/15}\leq(1-2/15)^{-1}=15/13$. The leading term $\ell=0$ in \eqref{eq:legendre-conv}, multiplied by $\sqrt{\sin\phi}$ and divided by $b_n$, therefore has sign $(-1)^k$ and magnitude greater than $49/50-1/6$, since the free cosine equals $(-1)^k\cos\eta$ and $\cos\eta\geq1-\eta^2/2>49/50$.

For $n\geq20$, the bounds on $c_j$ proved above give $b_n\geq 2/(3\sqrt n)$. For odd $n$, use $n/(n+1/2)\geq40/41$ and $c_{(n-1)/2}\geq1/\sqrt{2(n-1)}$. For $N\geq n/2$, they also give $b_N\leq2/\sqrt n$. Thus, in the same units, the sum of absolute values of terms $1\leq\ell\leq n/2$ is at most $3(7/6)\sum_{\ell\geq1}6^{-\ell}=7/10$. For $\ell>n/2$, the bound $|L_N(\cos\phi)|\leq1$ for all $N\geq0$ follows from the Laplace integral of $(\cos\phi+i\sin\phi\cos u)^N$ against $du/\pi$ on $[0,\pi]$~\cite[(18.10.5)]{dlmf}. This part has total at most
\[
(3/2)\sqrt n\sum_{\ell>n/2}6^{-\ell}
 \leq (9/5)n\,6^{-n/2}<1/100 .
\]
The last bound holds at $n=20$ and decreases thereafter. Together the two error bounds are smaller than $49/50-1/6$. Hence $G_n(-a)$ again has sign $(-1)^k$ at every node $\Psi_n(\phi)=k\pi$, $k\geq1$, in the stated phase range.
\end{proof}

\section{Counting roots and undoing the normalization}\label{sec:count}

For $n\geq20$, select nodes in increasing $\phi$ order, one for each $k\geq1$ with $k<2n/3$. This uses continuity and the limits \eqref{eq:phase-ends}. After any selected node for $k$, the intermediate value theorem gives a strictly later node for $k+1<2n/3$. Section~\ref{sec:branch} identifies increasing $\phi$ with increasing $a$. Lemma~\ref{lem:sign} gives alternating signs for $G_n(-a)$ at these nodes. The initial sign is positive at $a=0$. If $3\nmid n$, the last node is $k=D_n$, and this produces $D_n$ distinct negative roots. If $3\mid n$, the last node is $k=D_n-1$. The degree and its positive leading coefficient give the additional sign $(-1)^{D_n}$ at sufficiently large $a$, again producing $D_n$ distinct negative roots. The degree bound makes all the roots simple.

For \(0\le n\le19\), the finite verification in Section~\ref{sec:finite} gives the same result.

Define $T_n(q^k)=\binom nk q^k$ on real polynomials of degree at most $n$. Its homogeneous symbol $T_n[(z+w)^n]$, with $T_n$ acting on $z$, is $w^nJ_n(z/w)$, where
\[
J_n(q)=\sum_{k=0}^n\binom nk^2q^k=(1-q)^n L_n\big((1+q)/(1-q)\big).
\]
The Legendre zeros lie in $(-1,1)$~\cite[\S18.16(iii)]{dlmf}, so $J_n$ has only strictly negative zeros. The homogeneous symbol is a product of factors $z+bw$ with $b>0$, up to a positive constant, and is therefore real stable. Borcea and Br\"and\'en give real-rootedness preservation~\cite[Theorem~2(b), p.~470]{bb}. Thus $P_n=T_n(G_n)$ is real-rooted and has strictly negative zeros because its coefficients are nonnegative and its constant coefficient is $1$. To see simplicity, restrict $T_n$ to real polynomials of degree at most $D_n$, where the diagonal map is invertible. A neighborhood of the simple-rooted $G_n$ consists of real-rooted polynomials, as witnessed by sign changes in disjoint intervals about its roots. Therefore a neighborhood of $P_n$ consists of real-rooted polynomials. If $P_n$ had a repeated real factor $(q-\xi)^2$, replacing this factor by $(q-\xi)^2+\varepsilon$ for small $\varepsilon>0$ would contradict that neighborhood property.

\section{Exact finite verification}\label{sec:finite}

For \(0\le n\le19\), the checker constructs \(G_n\) by \eqref{eq:normalized-count} and the coefficients of \(P_n\) independently by adding one array row at a time. Let \(b_n(i_1,i_2,i_3)\) count arrays with the three indicated column sums. With \(b_0(\mathbf 0)=1\) and all other initial states zero, the transfer is
\[
b_{n+1}(\mathbf i)=\sum_{\substack{S\subseteq\{1,2,3\}\\|S|\le2}}
                 b_n(\mathbf i-\mathbf 1_S).
\]
States with negative coordinates are zero. All twenty diagonal vectors \(b_n(k,k,k)\) agree with \(\binom nk g_{n,k}\).

For each nonconstant polynomial, the verifier starts a Sturm chain with the polynomial and its derivative and repeatedly takes the negative remainder in \(\mathbb Q[q]\). A nonzero constant has a singleton chain. Every chain ends in a nonzero constant. For both polynomials the variation at \(-\infty\) is \(D_n\), the variation at \(0\) is zero, and the constant coefficient is positive. Thus every zero is simple and negative in these rows. Table~\ref{tab:prefix} shows the counts, with the endpoint signs and both coefficient vectors retained in the archive.

\begin{table}[htbp]
\centering
\small
\begin{tabular}{@{}r*{10}{r}@{}}
\toprule
\(n\) & 0&1&2&3&4&5&6&7&8&9\\
Distinct negative zeros &0&0&1&2&2&3&4&4&5&6\\
\midrule
\(n\) &10&11&12&13&14&15&16&17&18&19\\
Distinct negative zeros &6&7&8&8&9&10&10&11&12&12\\
\bottomrule
\end{tabular}
\caption{Exact root counts for both \(P_n\) and \(G_n\) below the analytic cutoff.}
\label{tab:prefix}
\end{table}

To reproduce the archive from the companion repository~\cite{code}, run the checker with standard Python 3.8 or later from the repository root.
\begin{verbatim}
python3 topics/balanced-arrays/code/verify_three_column_prefix.py
\end{verbatim}
The command emits one JSON record per row. Its output matches the repository archive \path{topics/balanced-arrays/results/three_column_prefix.jsonl}.

\paragraph{Independent check.}
An agent-assisted review verified the phase limits, cut integrals and their sign bounds, the degree count in each residue class, and the transfer of simplicity. It also checked both families by a distinct coefficient formula and rational Sturm chains through \(n=20\), including the first row of the analytic proof. The coefficient formula used for this check was
\[
a_{n,k}=\sum_{j=0}^{\lfloor n/2\rfloor}
   \binom n{2j}\binom{2j}{j}\binom{n-2j}{k-j}\binom{n-k+j}{k},
\]
where out-of-range binomial coefficients are zero. To derive it, set \(B=1+x_3\), \(x_1=\sqrt{u}\,w\), \(x_2=\sqrt{u}/w\) in the row polynomial. It becomes \(B+u+B\sqrt{u}(w+w^{-1})\). Only even powers of the last summand survive the constant term in \(w\). Choosing \(2j\) such factors and then the coefficient of \(u^k\) leaves \(B^{n-k+j}\), giving the required \(x_3^k\) coefficient.

Consecutive-row interlacing for \(U(2,3)\) and real-rootedness for all ranks and column counts remain open.

\begin{thebibliography}{9}
\bibitem{original} P.~Alexandersson, \emph{Real-rootedness of balanced arrays for uniform matroids}, RealRooted issue~2005, October 24, 2026. \href{https://github.com/PerAlexandersson/RealRooted/issues/2005}{Original conjecture}.
\bibitem{catalog} P.~Alexandersson, \emph{Real-rooted polynomials: open problems}, The Symmetric Functions Catalog, section on graphs, matroids and arrays. \href{https://www.symmetricfunctions.com/realRootedOpenProblems.htm\#openProblemsGraphs}{Conjecture entry}.
\bibitem{bb} J.~Borcea and P.~Br\"and\'en, P\'olya--Schur master theorems for circular domains and their boundaries, \emph{Ann. of Math.} (2) \textbf{170} (2009), no.~1, 465--492. \href{https://doi.org/10.4007/annals.2009.170.465}{doi:10.4007/annals.2009.170.465}.
\bibitem{dlmf} F.~W.~J. Olver et al.\ (eds.), \emph{NIST Digital Library of Mathematical Functions}, Chapters~15 and~18. \href{https://dlmf.nist.gov/}{DLMF}.
\bibitem{code} ProgramComputer, \emph{Balanced binary arrays}, Open Math Research, 2026. \href{https://github.com/ProgramComputer/open-math-research/pull/8/files}{Reviewed proof and exact computation}.
\end{thebibliography}
\end{document}
